% LAB BOREDOM — Section II appendices. GUIDED DRAFTING. Session 2026-08-21.
% Plan: plans/lab-boredom-appendices-plan-2026-08-21.md (all Phase-0 rulings received).
% ★ NUMBERING PROVISIONAL: these are SECTION-SCOPED designations (A.1/A.2/B/C/D, E optional);
%   the author will create other dissertation appendices later and the scheme WILL change
%   (memory: lab-boredom-appendix-renumbering). Every in-prose pointer stays ****** until
%   the dissertation-wide scheme exists — the designations below are working labels.
% Standing rulings: technical register permitted (D-G) · film-comparison tables EXCLUDED,
%   equipment-QC tables kept (D-B) · deployment-transfer topic DROPPED (D-C) · administered
%   survey wording is CANONICAL, no variant footnote (D-H) · papers bound as PDFs (D-E) ·
%   episode-only · data-as-agent · never "laboratory" · tree convention · no "ones".
% Tables: generated by analysis/appendix_tables.py -> out/appendix/D*.tex (spot-verified).
%
% BATCH A (headnote for A.1/A.2). REV 1 (2026-08-21, author note): "reanalyzed" ->
%   "reanalyzed and expanded upon."
% Verified facts: P1 authors King, Feeney, Tang, Das, Salvo (ASEE 2023 #37129); P2 authors
%   King, Lo, Das, Salvo (ASEE 2024 #44685, session-log-verified); IRB = UCI IRB Exempt
%   No. 20232678 (the 2024 form, per standing ruling); 3-of-8 pilot cohort fact; same
%   recordings reused in 2024; own publications in FIRST person.

Appendix A reproduces, in their published form, the two studies whose data this section
reanalyzed and expanded upon. Appendix A.1 is our 2023 pilot report, ``Work in Progress:
Physiological Assessment of Learning in a Virtual Reality Clinical Immersion Environment''
(King, Feeney, Tang, Das, and Salvo, ASEE 2023), which presented three of the first round's
eight participants. Appendix A.2 is our 2024 study, ``Assessment of Student Engagement in VR
Clinical Immersion Environments through Eye Tracking'' (King, Lo, Das, and Salvo, ASEE 2024),
which reported all twelve participants, drawing for the first round on the same recordings
the pilot had used. The analyses of this section pool those same twelve. Both studies were
conducted under UCI IRB Exempt No. 20232678.

% ============================================================================
% BATCH B (Appendix B — the survey instrument). Drafted 2026-08-21, awaiting author review.
% The seven questions are transcribed from the survey-response files and are IDENTICAL,
% word for word, across all twelve participants' records (cohort audit 2026-08-21, both
% rounds, one wording set). The seventh item's wording — open since L2 drafting — is hereby
% recovered and the ****** it carried resolves at integration.
% O-15 SLOT: the IRB-facing administration sentence is the author's, supplied at this
% batch's review — marked ****** below.
% Parsing counts verified against the tree 2026-08-21: exact 239 / converted 5 / flagged 1 /
% semantic 7 / missing 0 = 252. (Approved L2 prose still carries the pre-D-39 counts —
% flagged for a separate small revision batch with the author's L2 wording review.)

Appendix B reproduces the survey as it was administered. The seven questions below are
transcribed from the survey records themselves and are identical, word for word, across all
twelve participants' files. They were put to each participant aloud by the proctor and the
answers written down as given. The first question was asked just before each film began; the
remaining six were asked in the break that followed the film, their answers referring to the
film just watched.

\begin{enumerate}
\item How tired are you right now? (1-9; 1 completely awake, 9 completely exhausted)
\item On a scale from 1-9, how bored did you get watching this video? (1 not at all, 9
completely bored)
\item On a scale from 1-9, how interested and engaged in the video were you? (1 not engaged
at all, 9 intensely engaged)
\item Roughly how many minutes into the video until you started to get bored?
\item Did you fall asleep/fight to stay awake? (1 completely awake, 9 fell asleep)
\item Relatively speaking, roughly how long did the video feel to you?
\item On a scale from 1-9, how tired are you after watching the video (1 completely awake, 9
completely exhausted)
\end{enumerate}

Two quantities are derived rather than asked. Felt duration becomes a ratio against the
film's actual seventeen minutes, and depletion is the seventh answer minus the first. Because
the questions were read aloud and the answers recorded as free text, every answer passed
through a documented reading, and two coding rules govern the file. An answer of NA on any
1--9 item means the film produced none of the asked-after quality and is read as 1, the
scale's floor. A blank or NA on the minutes question means the participant never became
bored; those episodes are kept as their own category rather than assigned a number, since on
a duration item zero would mean bored from the first second. The reading of all 252 answers
(36 episodes by 7 items) breaks down as follows, and every interpreted value is stored
beside the exact words the participant gave:

\begin{tabular}{lcl}
\hline
reading & answers & meaning \\
\hline
exact & 239 & a plain number in range, no interpretation needed \\
converted & 5 & a unit or range converted (``25 sec'' to 0.42 minutes; ``10--12'' to 11) \\
flagged & 1 & an in-range answer carrying an anomaly, logged with its annotation \\
semantic & 7 & a meaningful non-number (never became bored; NA read as the scale floor) \\
missing & 0 & --- \\
\hline
\end{tabular}

% ============================================================================
% ★ APPENDIX B APPROVED by author 2026-08-21 (with the administration passage as revised
%   after the author's Q1-timing statement). O-15 RESOLVED. The four L2 revision edits
%   approved the same day.
%
% BATCH C1 (Appendix C — the criterion method in full). Drafted 2026-08-21, awaiting review.
% Base: PROSE-METHODS-CRITERION-VALIDITY-v2-2026-08-05.md §A ("do not rewrite" — numbers and
% logic stand). The file PREDATES the governing episode-only ruling (2026-08-07) and the
% D-39/D-40/D-41 re-runs, so the following DEVIATIONS are made and flagged to the author:
%   C1-a ¶1's opening two sentences (the films-apart question "answered by the preceding
%        tests") DELETED — those tests are excluded from the dissertation.
%   C1-b ¶5 closure figures refreshed to the current tree: +0.46/-0.42 -> +0.48/-0.44
%        (two-decimal style preserved; D-40 moved them).
%   C1-c ¶5's pupil sentence loses its film-separation clause -> "Pupil dilation does not
%        move with either report."
%   C1-d ¶6's closing clause recast episode-frame: "...and that no other recorded channel
%        does."
%   C1-e O-14 RESOLVED as ruled 2026-08-06: multiplicity handled in a closing FOOTNOTE, not
%        computed.
%   C1-f the v2 file's §C objections stay OUT (its own instruction: defense reserve, not
%        prose); §D notes are drafting apparatus, not appendix content.
% Everything else is v2 §A verbatim.
% BATCH C1 REV 1 (2026-08-21, author notes; C1-a..f all approved):
%   ¶2 pupil dropped from the between-person examples (not in any presented result) ->
%     closure share added; "participant-by-film cells" -> "episodes"; "three films" ->
%     "three episodes" here and in ¶6 (episode vocabulary).
%   ¶3 pupil example -> gaze/closure; "product-moment" glossed away ("Pearson's
%     correlation---the ordinary linear correlation coefficient"); NEW closing sentence on
%     what Spearman-Pearson agreement means (near-linearity: a 1->2 boredom step moves the
%     eye measures about as much as an 8->9 step).
%   ¶4 permutation mechanics restated as shuffling each participant's RATINGS among their
%     own episodes (was "film labels").
%   ¶5 "distortion largest where a reader would lean hardest" REPLACED with an explanatory
%     statement of the analytic-vs-permutation difference (textbook independence assumption
%     vs the data's own twelve-triplet structure).
%   NOTE: ¶5's pupil closing sentence ("Pupil dilation does not move with either report")
%     retained for now — it is the exclusion evidence; if the author wants pupil absent from
%     the appendix entirely, it and the D7b flat-channel table row go together (flagged).

Appendix C states the criterion method in full. The instruments were adopted because the
participants' own reports were treated as the criterion, and each channel was scored by
whether it agreed with what the participant said afterward. Restated in the form the
reprocessing can test: does a physiological channel move with a person's own account of their
boredom, episode by episode?

Both sides of the comparison are converted to within-subject standard scores before anything
is correlated. Each participant's value for one episode is expressed as its distance from
that participant's own three-episode mean, in units of that participant's own standard
deviation. The step is not cosmetic. Habitual gaze range, the share of an episode a person's
eyes ordinarily spend closed, and rating style vary enough between twelve people that an
uncentered correlation would chiefly measure who has a wide-ranging eye or a generous hand
with a nine-point scale. After centering, each participant contributes only the shape of
their three episodes relative to themselves, which is the quantity the question is about. The
association is then computed across all thirty-six episodes.

The reported coefficient is Spearman's rank correlation. It correlates the order of the values
rather than the values themselves, which suits this material for two reasons. No theory
predicts that self-reported boredom rises in proportion to degrees of gaze deviation or to
the share of an episode the eyes spend closed, so a statistic that assumes a straight line
assumes something the design does not license; and with twelve participants a single unusual
episode can bend a linear coefficient in a way it cannot bend a rank one. Pearson's
correlation---the ordinary linear correlation coefficient---is computed alongside and
reported in full in the tables that follow, both because the two coefficients should be
compared rather than one quietly chosen, and because their agreement is itself informative.
Across every channel and both criteria the median absolute difference between them is 0.032
and the largest is 0.117. Where a monotonic relationship is also close to linear, the rank
statistic loses almost nothing, so the choice of Spearman costs little. Nothing in it
flatters the result: on every association that carries the finding the Pearson coefficient is
the \textit{larger} of the two in magnitude, so the reported figure understates the
association rather than inflating it. The agreement between the two coefficients carries its
own meaning: the relationship is close to linear, so the change in the eye measures that
accompanies moving from a boredom rating of 1 to 2 is roughly the change that accompanies
moving from 8 to 9, rather than the association living only at the scale's extremes.

The probability attached to each coefficient is obtained by Monte Carlo permutation rather than from the
closed-form distribution, and the reason is a property of the design that the standard formula
cannot accommodate. Thirty-six paired observations are not thirty-six independent
observations: they are twelve participants contributing three films each. Within-subject
centering, which the comparison requires for the reasons given above, makes the dependence
structural rather than incidental---constraining three values to a mean of zero and a standard
deviation of one means that knowing two of them very nearly fixes the third. Both the rank and
the linear coefficient carry analytic p-values that assume independent observations, so both
are anti-conservative here. The correction is to build the null distribution from the data's
own structure: each participant's three ratings are shuffled among that participant's own
episodes, never across participants, the coefficient is recomputed, and the procedure repeats
twenty thousand times with a fixed seed. The proportion of shuffles producing an association
at least as strong as the observed one is the reported probability. Analytic values are
retained beside the Monte Carlo values throughout, so the size of the correction is visible
rather than absorbed.

That correction is consequential, and it is applied to both coefficients, which shows it to be
a feature of the design rather than an artifact of ranking. The two probabilities differ in
what they assume. The analytic value comes from the textbook formula, which treats the
thirty-six episodes as thirty-six independent observations; the Monte Carlo value is computed
from the data's own structure, with ratings reshuffled only inside each participant. Where an
association is weak the two barely differ; where it is strong, the analytic value runs three
to eight times too small, because the assumption of independence overstates how much evidence
thirty-six episodes from twelve people actually contain.
Three associations that clear the conventional threshold on the analytic value do not clear it
on the Monte Carlo value, and they are reported as not clearing it. What survives is stated
without hedging. On the per-participant baseline, the one assigned to this question, gaze
deviation moves with self-reported boredom at rho = +0.59 and against self-reported engagement
at rho = $-$0.59; the per-file baseline stands next at +0.52 and $-$0.53 and points the same
way, so the result does not turn on which of the two is used. Eye closure moves with reported
boredom at rho = +0.48 and against reported engagement at $-$0.44. All of these survive the
Monte Carlo test. Pupil dilation does not move with either report.

Two limits bound what any of this can be taken to show. The cohort is twelve, so these
coefficients describe this group and support no estimate of a population value; they are
reported for their direction and their relative ordering, not as parameters. And the
comparison is episode-level throughout---each physiological figure is an aggregate over a
whole seventeen-minute film and each rating is given in the break immediately afterward---so
nothing here time-locks a moment of experience to a moment of signal. What the analysis
establishes is that two of the eye channels track a person's own retrospective account across
their three episodes, and that no other recorded channel does.\footnote{The criterion tests are
reported as a family of descriptive associations, not as a set of independent hypothesis
tests, and the argument uses the ordering of channels rather than any individual probability.
No multiplicity correction is therefore applied to the family, by choice rather than
oversight.}

% ============================================================================
% ★ BATCH C1 APPROVED by author 2026-08-21 (rev 1; pupil sentence + appendix inclusion
%   KEPT, removable later at author discretion).
%
% BATCH C2 (Appendix C continued — the gaze-baseline statement). Drafted 2026-08-21,
% awaiting review. Base: PROSE-METHODS-GAZE-BASELINE-v2-2026-08-05.md §A ("do not rewrite").
% DEVIATIONS, flagged:
%   C2-a ¶3's separation assignment + ten-degree frequency result DELETED (film statistics,
%        excluded); device-forward recast as computed and reported WITH NO ASSIGNED QUESTION
%        (the architecture collapse recorded in the session log, stated openly).
%   C2-b ¶4 "strongest in the entire analysis" -> "strongest physiological association in
%        the analysis" (survey items run higher).
%   C2-c episode vocabulary (films -> episodes where the referent is the episode).
% Numbers verified current: +0.59/-0.59 per-participant, +0.52/-0.53 per-file, +0.36/-0.39
% device-forward (criterion-validity.csv; gaze channels unmoved by D-40).
% ★ STANDING TERMINOLOGY RULE (author, 2026-08-21, from a committee member): the
%   20,000-shuffle test is referred to as "MONTE CARLO" — rendered "Monte Carlo permutation"
%   at first technical naming, "Monte Carlo test/value" after. Applied across C1/C2 (6
%   substitutions). NOTE: the exact sign-flip test over 11 pairs is exhaustive, NOT Monte
%   Carlo — it keeps its name. Approved-prose sweep proposed to author: L1 x1, L2 x3, L6 x1.
% BATCH C2 REV 1 (2026-08-21, author notes):
%   P9 "and carrying it keeps the substitution ... visible" DELETED.
%   P10-1 opening sentence deleted; P10-2 the convention/one-name span rewritten in
%   explanatory form (the two concrete failure modes: the published-definition-only report
%   would have missed a real association; a bare "gaze deviation" number would leave the
%   reader unable to know which of three quantities it described).

Gaze deviation measures how far the eye has turned from center, and center is not something
the recording supplies. It has to be defined, and the definition is a modeling decision
rather than a detail of implementation. Three definitions are defensible here and all three
were computed. The first takes center to be the direction the headset itself is pointing,
which restores the published feature: the measure then asks how far the eye has turned away
from straight ahead. The second takes center to be the middle direction of that particular
recording, which is what the reprocessing initially computed. The third takes center to be
the middle direction of that participant across all three of their episodes, which preserves
however the headset happened to sit on that person's face while still exposing differences
between the episodes they watched.

The second definition carries a consequence worth stating plainly, because it is not obvious
and it is not small. When the reference is the recording's own middle direction, any posture
held for most of an episode disappears by construction. A participant who spends a whole
episode looking down and to the left has a median direction that is itself down and to the
left, and against that reference they register almost no deviation at all. The measure
therefore reports how much the eye \textit{moved around} within an episode and is blind to
where the eye was \textit{held}. The first definition has the opposite property: it retains
sustained posture, because its reference does not move with the participant. The third sits
between them, removing the posture a person holds across their whole session while keeping
the differences between their three episodes. None of the three is the correct measure of a
single underlying thing. They are measures of different things that one phrase---deviation
from center---had been concealing.

The assignment of measures to questions follows from how each is built rather than from how
each scored. The question this analysis asks of gaze is agreement: whether the eye's behavior
moves with what the participant reported about their own episode. That comparison is
within-subject throughout---each person is scored against themselves---and the participant's
own middle direction is the only one of the three references constructed the same way,
removing what differs between people while preserving what differs between one person's
episodes. The measure defined inside each single episode cannot serve that question, because
it removes the between-episode difference before the test can reach it. The device-forward
restoration is computed and reported in full beside the other two---it is the published
feature's own definition---but no question in this analysis is assigned to it.

Against the participant's own baseline the association with self-report is the strongest
physiological association in the analysis, at +0.59 against reported boredom and $-$0.59
against reported engagement, both surviving the Monte Carlo test comfortably; against the
recording's own middle direction it stands next at +0.52 and $-$0.53, pointing the same way,
so nothing in the finding depends on which of the two is chosen. Against the device's
forward axis---the published definition---the association does not survive the test, at
+0.36 and $-$0.39. These three outcomes are why the reporting rules matter. Had this
analysis reported only the published definition, a real association between gaze and the
participants' reports would have gone unseen. Had it reported any of the three under the
bare name gaze deviation, a reader could not have known which of three different quantities
the number described. Every gaze result in this section therefore names the baseline it was
computed on.

% ============================================================================
% BATCH C3 (Appendix C close — instrumentation bookkeeping). Drafted 2026-08-21, awaiting
% review. Content per D-B ruling (equipment QC kept): windows · rates/harmonization with the
% decimation measurement · closure definition + blink ceiling + its limits · poolability.
% Facts, all established: R1 window = middle ~13.6 min, ~2 min excluded inside each
% instructed closure (sync guard); R2 crop = whole film, instructed closures = stimulus
% boundaries (F-08); closure window-invariant; ~3 Hz vs 120 Hz (F-10, capture software not
% sensors); 5-point rolling median spans ~42 ms at R2's rate vs ~1.5 s at R1's; decimation
% factor 0.65–0.84 (gaze-baseline file §C, v5-verified); D-28 pooling rules incl. the shape
% caveat; both-eyes closure, R2 binary at source over 3,416,362 per-eye samples, R1 per-eye
% sentinel; 500 ms ceiling on VanderWerf 334±67 ms (watching video); R1 median interval
% ~334 ms -> blink rate unrecoverable in R1; blink-duration gaze-position dependency (E13).
% BATCH C3 REV 1 (2026-08-22, author notes):
%   P11-1 metatextual opener replaced — ¶11 opens on the three Round-1/Round-2 discrepancies.
%   P11-2 the instructed-closure PROCESS now described (L2 mentions the closures only as
%     brackets, checked line 404): proctor instruction at every boundary; film started on
%     "open your eyes", stopped on "close your eyes"; the closures mark episode boundaries.
%   P11-3 window-invariance anchored: every quantity computed twice (full window vs
%     R1-matched trim) and carried through the same tests; for closure the results are
%     identical throughout (window-comparison.csv: closure rows identical across windows —
%     anchored WITHOUT quoting the file's film-contrast statistics, per the episode ruling).
%   P13-1 "sentinel value" replaced with the concrete mechanism: a closed eye cannot be
%     read, and the tracker writes pupil diameter -1 in that eye's column, a value no real
%     eye can produce; closure read from that marker, both eyes at once.
% BATCH C3 REV 2 (2026-08-22):
%   P11 sentence replaced with the AUTHOR'S OWN WORDING (30-second closures; synchronization
%     with the screen recording for the gaze-depiction overlays) + one bridge sentence
%     restoring the boundary fact the window logic needs.
%   ★ P13 VERIFICATION (author-ordered, all 24 R1 cells): Round 1 carries BOTH a per-eye
%     open/close flag column (validL/validR, 0/1) AND the per-eye pupil columns reading -1.0
%     on closure. THE COMPUTATION USES ONLY THE -1 PUPIL MARKER (r1.py _closure: dilL==-1 &
%     dilR==-1; R1_CLOSED_SENTINEL=-1.0). EMPIRICAL CONTAINMENT: the flag NEVER marks both
%     eyes closed where the pupil columns do not (valid-only = 0 in all 24 cells); the pupil
%     marker catches 2,381 additional both-closed samples of 25,187 total (~9%; worst-case
%     cell agreement 87.8%, S07 BOR). The -1 definition is a strict superset — nothing the
%     dedicated flag calls closure is dropped. Findings reported to author; ¶13 revision
%     PROPOSED, awaiting decision.
% BATCH C3 REV 3 (2026-08-22, author ruling on the P13 findings): ¶13 describes ONLY the -1
%   pupil marker as the data from which Round 1 closure was computed — all mention of the
%   dedicated open/close flag column REMOVED from prose (the proposed two-record revision is
%   discarded; the rev-1 pupil-marker-only text STANDS). The verification record above stays
%   in this ledger as the audit trail for the choice.
% ★ ¶¶11–13 APPROVED 2026-08-22. ★★ APPENDIX C COMPLETE — 13 ¶¶.
%
% ============================================================================
% BATCH D-i (Appendix D framing + tables D1–D4). Drafted 2026-08-22, awaiting review.
% ★ STANDING RULE (author, 2026-08-22): EVERY table introduction must be articulated so a
%   reader with NO statistical training understands what each category is doing and what the
%   results mean — keys for every column label, plain statements of what the numbers mean.
% Film codes defined in the framing per the author's wording (BOR = 1989 Microsoft Word
% tutorial; CLC = 3D VR spinal deformation surgery; INT = government-produced alien
% reproduction vehicles video).
% Tables generated by analysis/appendix_tables.py (D1 labels de-jargoned + % hazard fixed;
% D7a columns renamed and pruned; D3 paired tally flipped to predicted-direction
% agree/contra/tied). REV of 2026-08-22: D-i intros rewritten to the lay standard.

Appendix D collects the full numerical tables behind the findings. Three film codes recur as
episode identifiers and carry no analytical weight: BOR is the boring film sample, a 1989
Microsoft Word tutorial; CLC is the clinical film sample, 3D VR footage of spinal deformation
surgery; INT is the interesting film sample, a video about government-produced alien
reproduction vehicles. Unless a column states otherwise, every table covers all thirty-six
episodes from all twelve participants. Two statistical conventions hold throughout.
Correlations are computed on each participant's own within-person scale, the standardization
Appendix C describes; and each correlation carries a Monte Carlo probability, written p: the
share of twenty thousand chance reshuffles of the data that produced an association at least
as strong as the observed one, so a small p means chance is an unlikely explanation. The
exact sign test's probabilities are exhaustive counts rather than samples.

Table D1 reports, for every recorded channel, how closely that channel moved with what
participants reported. Two correlation statistics appear side by side: $\rho$ (rho) is
Spearman's rank correlation and $r$ is Pearson's linear correlation. Both run from $-$1 to
+1, where +1 means the channel rose and fell in perfect step with the rating, $-$1 means it
moved perfectly opposite, and 0 means the two were unrelated. The subscript names the rating:
bor is reported boredom, eng is reported engagement. Each p is the Monte Carlo probability
attached to the $\rho$ beside it---so p = 0.0024 means about one chance reshuffle in four
hundred matched the observed association. The column $n$ is the number of episodes with
usable values for that channel. Read as a whole, the table shows the two findings the body
reports: gaze deviation against the participant's own baseline and eye closure move with both
ratings at probabilities chance cannot comfortably explain, and no other channel does.

% % D1-criterion-full — every channel against both ratings; Spearman with within-participant Monte Carlo permutation p, Pearson beside it; pooled 36 episodes
% generated by analysis/appendix_tables.py, seeds fixed, tree convention
\begin{tabular}{lccccccc}
\hline
channel & $n$ & $\rho_{bor}$ & $p$ & $r_{bor}$ & $\rho_{eng}$ & $p$ & $r_{eng}$ \\
\hline
pupil dilation & 36 & -0.172 & 0.4037 & -0.164 & +0.036 & 0.8608 & +0.100 \\
eye closure & 36 & +0.475 & 0.0173 & +0.580 & -0.435 & 0.0279 & -0.535 \\
gaze deviation (per-file baseline) & 36 & +0.524 & 0.0074 & +0.574 & -0.527 & 0.0075 & -0.590 \\
gaze deviation (device-forward baseline) & 36 & +0.359 & 0.0751 & +0.341 & -0.391 & 0.0515 & -0.403 \\
gaze deviation, 90th percentile (device-forward) & 36 & -0.046 & 0.8222 & +0.000 & -0.033 & 0.8730 & -0.088 \\
share of episode beyond 10 degrees (device-forward) & 36 & +0.324 & 0.1114 & +0.346 & -0.334 & 0.0990 & -0.389 \\
excursions beyond 10 degrees per minute (device-forward) & 36 & -0.082 & 0.6829 & -0.148 & +0.184 & 0.3622 & +0.216 \\
gaze deviation variance (per-file) & 36 & -0.114 & 0.5675 & -0.105 & +0.106 & 0.6073 & +0.091 \\
cognitive load & 36 & -0.092 & 0.6516 & -0.059 & +0.107 & 0.5954 & +0.097 \\
heart rate & 34 & +0.088 & 0.6774 & +0.079 & +0.036 & 0.8647 & +0.052 \\
gaze deviation (per-participant baseline) & 36 & +0.590 & 0.0024 & +0.623 & -0.585 & 0.0022 & -0.609 \\
\hline
\end{tabular}


Table D2 splits the two load-bearing correlations by recording round. Each cell holds the
correlation with its Monte Carlo probability in parentheses. The point of the split is
visible across each row: both rounds move in the same direction as the pooled result, so the
pooled finding is the property of neither round alone---and Round 2, recorded at the sensor's
full rate, agrees most strongly.

% % D2-per-round — the two load-bearing channels split by round; Round 2 significance is per the 12-episode within-participant permutation; no raw magnitudes cross rounds
% generated by analysis/appendix_tables.py, seeds fixed, tree convention
\begin{tabular}{llccc}
\hline
channel & rating & Round 1 ($n$=24) & Round 2 ($n$=12) & pooled ($n$=36) \\
\hline
gaze deviation (per-participant baseline) & boredom & +0.475 (0.0552) & +0.868 (0.0049) & +0.590 (0.0024) \\
gaze deviation (per-participant baseline) & engagement & -0.429 (0.0839) & -0.928 (0.0067) & -0.585 (0.0022) \\
eye closure & boredom & +0.420 (0.0916) & +0.538 (0.1261) & +0.475 (0.0173) \\
eye closure & engagement & -0.397 (0.1105) & -0.452 (0.2231) & -0.435 (0.0279) \\
\hline
\end{tabular}


Table D3 asks of every survey item what Table D1 asks of every channel: did the item move
with the two ratings, and with the two eye measures? The first four columns are correlations
with Monte Carlo probabilities, read exactly as in Table D1. The last two columns report a
simpler test. For each participant, take the episode that participant rated most boring and
the episode that participant rated most engaging, and ask whether the item differed between
them in the direction boredom predicts---more sleep fought, more fatigue afterward, boredom
arriving sooner. The paired column counts participants: agree means the item moved the
predicted way, contra means it moved the opposite way, tied means the participant gave
identical values in both episodes. The sign p is the chance of a split at least that lopsided
if the direction were decided by coin flip: eleven of twelve going one way, as the sleep item
shows, happens by chance about once in a thousand.

% % D3-survey-battery — every survey item against both ratings and both eye measures (rho, Monte Carlo permutation p), with the within-participant paired tally in each item's predicted direction (agree/contra/tied) and exact sign p
% generated by analysis/appendix_tables.py, seeds fixed, tree convention
\begin{tabular}{lcccccc}
\hline
item & vs boredom & vs engagement & vs gaze & vs closure & paired & sign $p$ \\
\hline
sleep fight & +0.770 ($<$0.0001) & -0.749 ($<$0.0001) & +0.461 (0.0192) & +0.665 (0.0006) & 11/0/1 & 0.0010 \\
fatigue after & +0.643 (0.0008) & -0.615 (0.0018) & +0.420 (0.0348) & +0.690 (0.0006) & 10/1/1 & 0.0117 \\
minutes until bored & -0.778 (0.0001) & +0.744 (0.0006) & -0.375 (0.0895) & -0.529 (0.0150) & 8/0/0 & 0.0078 \\
depletion & +0.521 (0.0120) & -0.443 (0.0329) & +0.266 (0.2077) & +0.418 (0.0448) & 9/2/1 & 0.0654 \\
fatigue prior & +0.069 (0.7411) & -0.207 (0.2975) & +0.320 (0.1071) & +0.238 (0.2399) & 6/3/3 & 0.5078 \\
felt duration ratio & +0.494 (0.0107) & -0.550 (0.0058) & +0.241 (0.2344) & +0.092 (0.6491) & 9/2/1 & 0.0654 \\
\hline
\end{tabular}


Table D4 reports the section's central sorting of the episodes. For each episode, the two eye
measures were combined into a single reading of whether the body's watching looked engaged or
bored relative to that participant's own average, and that reading was set beside what the
participant's ratings said. Four outcomes are possible, and the top half of the table counts
them: concordant-bored means the eyes and the ratings both say bored; concordant-engaged
means both say engaged; divergent means the eyes read engaged while the participant reported
boredom; reverse means the eyes read bored while the participant reported no boredom. The
bottom half answers a fairness question about the combination: if the sorting is repeated
using eye closure by itself, or gaze deviation by itself, does the divergent category
survive? It does---twelve episodes under closure alone, eight under gaze alone, nine under
the two combined---so the category is not manufactured by the choice of how the two measures
were weighted.

% % D4-decomposition-sweep — the four-category decomposition across all 36 episodes, and the divergent count under each single-measure definition (weighting robustness)
% generated by analysis/appendix_tables.py, seeds fixed, tree convention
\begin{tabular}{lc}
\hline
category / definition & episodes \\
\hline
concordant-bored & 13 \\
concordant-engaged & 11 \\
divergent & 9 \\
reverse & 3 \\
divergent under eye closure alone & 12 \\
divergent under gaze deviation alone & 8 \\
divergent under both (the composite) & 9 \\
\hline
\end{tabular}


% ============================================================================
% ★ BATCH D-i APPROVED by author 2026-08-22 (framing + D1–D4 with lay introductions).
%
% BATCH D-ii (tables D5–D8 with lay introductions per the standing rule). Drafted
% 2026-08-22, awaiting review. D5's minutes column formatting fixed (0.416667 -> 0.42, the
% parsing convention). D7a retention range verified from the tree: 36.3–74.4%, 3 of 24
% recordings below 50%.
% BATCH D-ii REV 1 (2026-08-22, author notes): D6 opening rephrased — no
%   "answers-a-question-a-reader-will-ask" framing; states plainly what the table
%   illustrates. ALL ELSE APPROVED.
% ★ AUTHOR ADVISORY (2026-08-22): the EEG material (D7a, possibly D7b's EEG-adjacent role)
%   will LIKELY BE REMOVED later — the published papers already state why EEG was dropped,
%   and the dissertation need not rehash it. LEFT IN for now by author order; expect a
%   removal decision at a later pass.
% ★★ APPENDICES A.1/A.2, B, C, D ALL DRAFTED AND APPROVED 2026-08-22 (E remains optional,
%   not drafted). Remaining: integration (designation pins stay ****** per the renumbering
%   rule; L1 re-render; assembled-document regeneration; commit gate).

Table D5 lists all thirty-six episodes, one row each, so any claim in the body can be
checked against the raw picture of a single episode. The columns give: the participant's
code (subj); the film code (ep); clo\%, the percentage of the episode during which both eyes
were closed; gaze, the median distance of the eye from that participant's own habitual
direction, in degrees; the participant's ratings of boredom (bor) and engagement (eng), each
on the 1--9 scale; slp, how hard the participant fought sleep, 1--9; min, the reported
minutes until boredom arrived, with -- meaning it never did; felt, the felt duration as a
fraction of the film's actual seventeen minutes, so 1.18 means the film felt about twenty
minutes long and 0.47 means it felt like about eight; fat$_0$ and fat$_1$, tiredness just
before and just after the episode, 1--9; and the episode's category from Table D4's sorting.

% % D5-per-participant — all 36 episodes: eye measures (closure %, gaze deviation deg), all seven items (felt duration as the ratio), and the decomposition category; film code is an episode identifier only
% generated by analysis/appendix_tables.py, seeds fixed, tree convention
\begin{tabular}{llcccccccccl}
\hline
subj & ep & clo\% & gaze & bor & eng & slp & min & felt & fat$_0$ & fat$_1$ & category \\
\hline
R2-01 & BOR & 18.4 & 10.7 & 8 & 2 & 7 & 4 & 1.18 & 3 & 5 & concordant-bored \\
R2-01 & CLC & 3.1 & 10.0 & 6 & 4 & 5 & 7 & 1.18 & 2 & 3 & divergent \\
R2-01 & INT & 8.3 & 8.1 & 3 & 6 & 2 & 13 & 1.06 & 2 & 3 & concordant-engaged \\
R2-02 & BOR & 24.7 & 11.2 & 9 & 1 & 9 & 0.5 & 1.76 & 6 & 9 & concordant-bored \\
R2-02 & CLC & 3.6 & 8.3 & 8 & 4 & 7 & 7 & 1.76 & 6 & 8 & divergent \\
R2-02 & INT & 2.3 & 8.5 & 7 & 6 & 6 & 5 & 1.18 & 5 & 7 & divergent \\
R2-03 & BOR & 11.7 & 10.6 & 3 & 6 & 3 & -- & 1.18 & 5 & 6 & reverse \\
R2-03 & CLC & 6.4 & 6.9 & 1 & 8 & 1 & -- & 1.18 & 4 & 5 & concordant-engaged \\
R2-03 & INT & 5.5 & 7.5 & 2 & 8 & 1 & -- & 0.88 & 3 & 4 & concordant-engaged \\
R2-04 & BOR & 7.8 & 7.8 & 7 & 1 & 6 & 2 & 1.18 & 3 & 5 & concordant-bored \\
R2-04 & CLC & 5.5 & 6.8 & 2 & 3 & 1 & 7 & 0.88 & 3 & 3 & concordant-engaged \\
R2-04 & INT & 7.1 & 6.3 & 3 & 7 & 4 & 8.5 & 0.88 & 2 & 5 & concordant-engaged \\
S01 & BOR & 9.0 & 10.7 & 6 & 3 & 2 & 5 & 0.88 & 5 & 5 & concordant-bored \\
S01 & CLC & 4.3 & 8.3 & 3 & 5 & 2 & 10 & 1.18 & 4 & 5 & concordant-engaged \\
S01 & INT & 7.0 & 8.1 & 2 & 7 & 1 & -- & 0.59 & 5 & 4 & concordant-engaged \\
S02 & BOR & 17.9 & 10.5 & 9 & 2 & 7 & 4 & 1.18 & 6 & 7 & concordant-bored \\
S02 & CLC & 19.0 & 10.8 & 8 & 3 & 5 & 3 & 1.76 & 4 & 6 & concordant-bored \\
S02 & INT & 5.2 & 9.3 & 3 & 6 & 1 & -- & 1.76 & 4 & 4 & concordant-engaged \\
S03 & BOR & 15.9 & 11.2 & 9 & 1 & 8 & 0.416667 & 3.53 & 3 & 8 & concordant-bored \\
S03 & CLC & 5.4 & 11.0 & 3 & 7 & 2 & 8 & 0.88 & 6 & 3 & concordant-engaged \\
S03 & INT & 12.2 & 9.8 & 3 & 7 & 7 & 2 & 1.76 & 5 & 7 & concordant-engaged \\
S04 & BOR & 8.2 & 12.7 & 1 & 7 & 1 & -- & 0.47 & 7 & 1 & reverse \\
S04 & CLC & 4.8 & 16.0 & 9 & 1 & 6 & 1 & 1.76 & 7 & 9 & concordant-bored \\
S04 & INT & 5.0 & 11.3 & 3 & 7 & 1 & 10 & 1.18 & 4 & 5 & concordant-engaged \\
S05 & BOR & 77.8 & 9.5 & 9 & 2 & 9 & 3 & 0.29 & 3 & 8 & concordant-bored \\
S05 & CLC & 11.6 & 12.9 & 6 & 5 & 1 & 5 & 1.00 & 3 & 5 & concordant-bored \\
S05 & INT & 13.9 & 8.8 & 5 & 4 & 3 & 7 & 1.18 & 6 & 5 & divergent \\
S06 & BOR & 15.7 & 17.5 & 4 & 5 & 6 & 5 & 1.18 & 8 & 8 & reverse \\
S06 & CLC & 7.0 & 18.2 & 6 & 3 & 3 & 4 & 1.47 & 6 & 5 & divergent \\
S06 & INT & 10.1 & 9.8 & 8 & 2 & 6 & 1 & 1.76 & 4 & 6 & divergent \\
S07 & BOR & 46.6 & 11.7 & 9 & 2 & 9 & 2 & 1.76 & 3 & 7 & concordant-bored \\
S07 & CLC & 10.6 & 7.0 & 6 & 7 & 3 & 5 & 1.18 & 4 & 3 & divergent \\
S07 & INT & 21.4 & 10.6 & 7 & 6 & 6 & 10 & 0.65 & 4 & 3 & concordant-bored \\
S08 & BOR & 57.7 & 13.5 & 9 & 1 & 8 & 0.5 & 1.47 & 7 & 8 & concordant-bored \\
S08 & CLC & 11.3 & 7.5 & 7 & 3 & 7 & 2 & 1.18 & 6 & 7 & divergent \\
S08 & INT & 10.5 & 10.3 & 5 & 4 & 5 & 6 & 0.88 & 4 & 6 & divergent \\
\hline
\end{tabular}


Table D6 illustrates how much the results depend on S06, the one participant whose
measurements contradicted S06's own reports throughout. Every
load-bearing correlation is computed twice, once with all twelve participants and once with
S06 removed and the remaining eleven re-standardized from scratch. Read down the two
columns: every association is stronger without S06, and none changes direction or loses
significance with S06 included. Retaining S06, which this analysis does throughout,
therefore makes every reported figure the more conservative of the two. The bottom four
rows repeat the check for the paired comparison of Table D3, with the same result.

% % D6-s06-comparison — every load-bearing association with S06 included (n=12) and excluded (n=11, re-standardized from scratch); S06 is retained in all reported results
% generated by analysis/appendix_tables.py, seeds fixed, tree convention
\begin{tabular}{llcc}
\hline
measure & rating & with S06 & without S06 \\
\hline
gaze deviation & boredom & +0.590 (0.0022) & +0.701 (0.0006) \\
gaze deviation & engagement & -0.585 (0.0026) & -0.682 (0.0008) \\
eye closure & boredom & +0.475 (0.0181) & +0.552 (0.0065) \\
eye closure & engagement & -0.435 (0.0290) & -0.512 (0.0107) \\
sleep-fight & boredom & +0.770 ($<$0.0001) & +0.857 ($<$0.0001) \\
sleep-fight & engagement & -0.749 (0.0001) & -0.851 ($<$0.0001) \\
fatigue after & boredom & +0.643 (0.0003) & +0.771 (0.0001) \\
fatigue after & engagement & -0.615 (0.0012) & -0.749 (0.0001) \\
gaze deviation & paired tally & 10/2/0 (0.0386) & 10/1/0 (0.0117) \\
eye closure & paired tally & 10/2/0 (0.0386) & 10/1/0 (0.0117) \\
sleep-fight & paired tally & 11/0/1 (0.0010) & 11/0/0 (0.0010) \\
fatigue after & paired tally & 10/1/1 (0.0117) & 10/0/1 (0.0020) \\
\hline
\end{tabular}


Table D7a documents the EEG's data quality, which is the evidence behind the decision to
treat that channel as this analysis treats it. Each recording was cut into short epochs, and
an epoch contaminated by artifact---movement, muscle activity, a poorly seated
electrode---was rejected. The columns give the number of epochs in each recording, the
number retained after rejection, the percentage retained, and the ratio of artifact variance
to clean-signal variance, where a larger number means a noisier recording. Retention runs
from 36.3\% to 74.4\% across the twenty-four recordings, with three below half. No recording
was excluded for noise, because these figures are themselves the finding: the channel's
thinness is documented rather than hidden. Table D7b then prints the criterion correlations
for the three excluded channels, read exactly as Table D1; the numbers are the reason the
exclusions are not arbitrary, since none of the three moves with either rating at anything
beyond chance.

% % D7a-eeg-retention — EEG epoch retention by recording (the evidence behind D-16); ratio = artifact variance over clean variance
% generated by analysis/appendix_tables.py, seeds fixed, tree convention
\begin{tabular}{llcccc}
\hline
subject & episode & epochs & retained & retained \% & artifact/clean variance \\
\hline
S01 & BOR & 534 & 336 & 62.9 & 6.6 \\
S01 & CLC & 539 & 305 & 56.6 & 15.8 \\
S01 & INT & 502 & 273 & 54.4 & 7.8 \\
S02 & BOR & 513 & 262 & 51.1 & 4.9 \\
S02 & CLC & 508 & 321 & 63.2 & 3.4 \\
S02 & INT & 510 & 289 & 56.7 & 7.5 \\
S03 & BOR & 523 & 388 & 74.2 & 12.2 \\
S03 & CLC & 482 & 334 & 69.3 & 11.2 \\
S03 & INT & 481 & 358 & 74.4 & 9.7 \\
S04 & BOR & 510 & 321 & 62.9 & 2.9 \\
S04 & CLC & 514 & 319 & 62.1 & 3.9 \\
S04 & INT & 507 & 357 & 70.4 & 3.5 \\
S05 & BOR & 517 & 340 & 65.8 & 83.3 \\
S05 & CLC & 482 & 224 & 46.5 & 7.5 \\
S05 & INT & 505 & 345 & 68.3 & 60.3 \\
S06 & BOR & 510 & 324 & 63.5 & 6.0 \\
S06 & CLC & 505 & 323 & 64.0 & 6.0 \\
S06 & INT & 569 & 402 & 70.7 & 10.9 \\
S07 & BOR & 568 & 418 & 73.6 & 4.4 \\
S07 & CLC & 558 & 398 & 71.3 & 7.5 \\
S07 & INT & 510 & 347 & 68.0 & 28.5 \\
S08 & BOR & 509 & 336 & 66.0 & 20.5 \\
S08 & CLC & 520 & 189 & 36.3 & 12.0 \\
S08 & INT & 504 & 222 & 44.0 & 7.6 \\
\hline
\end{tabular}

% % D7b-flat-channels — the excluded channels' criterion rows — the numbers that make the exclusions non-arbitrary
% generated by analysis/appendix_tables.py, seeds fixed, tree convention
\begin{tabular}{lcc}
\hline
channel & vs boredom & vs engagement \\
\hline
pupil dilation & -0.172 (0.4037) & +0.036 (0.8608) \\
cognitive load & -0.092 (0.6516) & +0.107 (0.5954) \\
HR & +0.088 (0.6774) & +0.036 (0.8647) \\
\hline
\end{tabular}


Table D8 counts closure events per round. Every interval during which both eyes were shut is
one event; events at or below the 500-millisecond ceiling are blink-length, longer events
are extended closures, and the final column gives the round's longest single event in
seconds. The two rows are printed separately and never summed, because the counts are not
comparable across rounds: Round 1's slow sampling means a brief closure there spans at most
a sample or two, so Round 1 undercounts short events by construction.

% % D8-closure-episodes — closure events per round: total, blink-length (≤500 ms), extended, longest single event (s)
% generated by analysis/appendix_tables.py, seeds fixed, tree convention
\begin{tabular}{lcccc}
\hline
round & events & blink-length & extended & longest (s) \\
\hline
Round 1 & 4365 & 3501 & 864 & 449.2 \\
Round 2 & 5131 & 4895 & 236 & 29.7 \\
\hline
\end{tabular}


Round 1 and Round 2 differ in three respects that bear on the numbers: the recording window,
the sampling rate, and the way eye closure reaches the file. The window difference begins
with the protocol's own markers. At the very start and end of each video sample, the proctor
asked the subject to close their eyes for 30 seconds. This was done to help synchronize the
data collection with the screen recording of the video session so that we could create
accurate gaze depiction overlays of where the students were looking while watching each
respective video. Those instructed closures mark each episode's boundaries in the recordings.
Round 2's recordings are cropped to those boundaries and cover essentially the whole film. Round 1's recordings carry a synchronization guard adopted while
the protocol was still being settled: roughly two minutes are excluded inside each instructed
closure, which leaves the middle thirteen and a half minutes or so of each seventeen-minute
film. A guard of that size could matter for a measure that accumulates, so every quantity was
computed twice---once on the full Round 2 recording and again on a window trimmed to match
Round 1's---and carried through the same tests both ways. For eye closure the two
computations give the same results throughout: no figure this section reports moves when the
window changes, which is what licenses treating the two rounds' closure measurements as one
family.

The second concerns sampling rates. Round 1's telemetry logs at roughly 3 Hz and Round 2's at
120 Hz, the sensor's own rate; the difference belongs to the capture software, not to the
sensors. Three rules follow. Means of per-sample values pool across rounds; variances and
windowed statistics do not. The gaze feature is a windowed statistic---its five-point rolling
median spans about forty-two milliseconds at Round 2's rate and about a second and a half at
Round 1's---and recomputing the identical feature on Round 2 data decimated to Round 1's rate
moves its value by a factor of 0.65 to 0.84, so no raw gaze magnitude crosses rounds anywhere
in this analysis. The criterion analysis pools the rounds legitimately, because
within-subject standardization is scale-invariant and removes a multiplicative rate
difference by construction; the stated limit is that this does not guarantee removal of a
rate effect on the shape of a participant's three-episode pattern, and the per-round split in
the tables that follow is what bounds that concern.

The third concerns eye closure. Closure means both eyes shut at once, on both rounds. Round 2
records closure binary at source---0 or 1 across 3,416,362 per-eye samples---so no threshold
was tuned. Round 1's tracker marks a closed eye differently: when an eye is shut the tracker
cannot read it, and it writes a pupil diameter of $-$1 in that eye's column, a value no real
eye can produce. Closure in Round 1 is read from that marker, both eyes at once. A ceiling of
500 milliseconds separates blink-length events from extended closures, licensed by
VanderWerf's measurement of spontaneous blinks at 334 plus-or-minus 67 milliseconds in
subjects watching video; every event at or below the ceiling counts as a blink, everything
longer as an extended closure. Two limits attach. Round 1's median sampling interval is
itself about the length of a blink, so a blink occupies zero or one samples there, and blink
rate is unrecoverable in Round 1 at any threshold---only extended closures are. The ceiling
is a duration threshold, and blink duration varies systematically with vertical gaze
position, which is controllable at a fixed screen but not in a headset where gaze is the
participant's own behavior; the blink-against-extended partition and the event counts in the
tables that follow inherit that dependency.
